\PassOptionsToPackage{table}{xcolor}
\documentclass[sigconf,
natbib=true]{acmart}

\AtBeginDocument{%
  \providecommand\BibTeX{{%
    \normalfont B\kern-0.5em{\scshape i\kern-0.25em b}\kern-0.8em\TeX}}}

\usepackage{enumitem}
\usepackage{csquotes}
\usepackage{tikz}
\usepackage{pifont}
\usepackage{amsmath}
\usepackage{calc}
\usepackage{booktabs}
\usepackage{ifthen}
\usepackage{wrapfig}
\usepackage{pgfplots}
\usepackage{pgfplotstable}
\usepackage{numprint}

\usetikzlibrary{patterns}
\usetikzlibrary{matrix}
\usetikzlibrary{arrows}
\usepgfplotslibrary{statistics}

\usetikzlibrary{calc,fit,backgrounds,positioning,shapes,arrows.meta,shapes.multipart}
\usepgfplotslibrary{dateplot}

\makeatletter
\pgfplotsset{
	boxplot/hide outliers/.code={
		\def\pgfplotsplothandlerboxplot@outlier{}%
	}
}
\tikzset{nomorepostaction/.code=\let\tikz@postactions\pgfutil@empty}
\tikzset{nomorepreaction/.code=\let\tikz@preactions\pgfutil@empty}
\makeatother

\usepackage{multirow}
\usepackage{makecell}

\newcommand*\rot{\rotatebox{90}}
\newcommand*\rotc{\rotatebox[origin=rb]{270}}
\newcommand*\OK{$\bullet$}
\newcommand*\UP{\ding{218}}
\newcommand*\DOWN{\ding{216}}
\newcommand*\rhombusOK{\ding{70}}
\newcommand*\significOK{$\star$}
\newcommand*\notsignificOK{$\diamond$}

\newcommand*\PARTIAL{$\circ$}
\newcommand*\NA{\cellcolor{gray!40}}
\newcommand*\UNKNOWN{--}

\definecolor{tablelightgray}{gray}{0.8}
\newcommand*\tls{\color{tablelightgray}\vrule width 0.4pt} %

\newcommand{\mc}[2]{\multicolumn{#1}{c|}{\makecell{#2}}} %
\newcommand{\mct}[2]{\mc{#1}{\tiny #2}} %
\newcommand{\rott}[1]{\rot{\tiny #1}}

\newlength{\largestheaderlength}%

\definecolor{applegreen}{rgb}{0.55, 0.71, 0.0}

\newcommand{\addcite}{~\cite{?}}
\renewcommand{\addcite}{}

\pgfplotsset{compat=1.16}

\graphicspath{{figures/}}

\copyrightyear{2020}
\acmYear{2020}
\setcopyright{rightsretained}
\acmConference[VINCI 2020]{The 13th International Symposium on Visual
Information Communication and Interaction}{December 8--10,
2020}{Eindhoven, Netherlands}
\acmBooktitle{The 13th International Symposium on Visual Information
Communication and Interaction (VINCI 2020), December 8--10, 2020,
Eindhoven, Netherlands}\acmDOI{10.1145/3430036.3430054}
\acmISBN{978-1-4503-8750-7/20/12}

\begin{document}

\title{Survey on User Studies on the Effectiveness of Treemaps}

\author{Carolin Fiedler}
\email{mail@carolinfiedler.de}
\affiliation{%
  \institution{Hasso Plattner Institute,\\Faculty of Digital Engineering,\\University of Potsdam,\\Germany}
}

\author{Willy Scheibel}
\email{willy.scheibel@hpi.de}
\affiliation{%
  \institution{Hasso Plattner Institute,\\Faculty of Digital Engineering,\\University of Potsdam,\\Germany}
}

\author{Daniel Limberger}
\email{daniel.limberger@hpi.de}
\affiliation{%
  \institution{Hasso Plattner Institute,\\Faculty of Digital Engineering,\\University of Potsdam,\\Germany}
}

\author{Matthias Trapp}
\email{matthias.trapp@hpi.de}
\affiliation{%
  \institution{Hasso Plattner Institute,\\Faculty of Digital Engineering,\\University of Potsdam,\\Germany}
}

\author{J{\"u}rgen D{\"o}llner}
\email{juergen.doellner@hpi.de}
\affiliation{%
  \institution{Hasso Plattner Institute,\\Faculty of Digital Engineering,\\University of Potsdam,\\Germany}
}

\renewcommand{\shortauthors}{C. Fiedler, W. Scheibel, D. Limberger, M. Trapp, and J. D{\"o}llner}

\begin{abstract}
    Treemaps are a commonly used tool for the visual display and communication of tree-structured, multi-variate data.
 	In order to confidently know when and how treemaps can best be applied, the research community uses usability studies and controlled experiments to \enquote{understand the potential and limitations of our tools}~(Plaisant, 2004).
    To support the communities' understanding and usage of treemaps, this survey provides a comprehensive review and detailed overview of 69 user studies related to treemaps.
    However, due to pitfalls and shortcomings in design, conduct, and reporting of the user studies, there is little that can be reliably derived or accepted as a generalized statement.
    Fundamental open questions include configuration, compatible tasks, use cases, and perceptional characteristics of treemaps.
    The reliability of findings and statements is discussed and common pitfalls of treemap user studies are identified.
\end{abstract}

\begin{CCSXML}
    <ccs2012>
    <concept>
    <concept_id>10003120.10003145.10003146.10010889</concept_id>
    <concept_desc>Human-centered computing~Treemaps</concept_desc>
    <concept_significance>500</concept_significance>
    </concept>
    <concept>
    <concept_id>10003120.10003145.10011770</concept_id>
    <concept_desc>Human-centered computing~Visualization design and evaluation methods</concept_desc>
    <concept_significance>500</concept_significance>
    </concept>
    </ccs2012>
\end{CCSXML}

\ccsdesc[500]{Human-centered computing~Treemaps}
\ccsdesc[500]{Human-centered computing~Visualization design and evaluation methods}

\keywords{Treemap, Evaluation, Literature Review}

\maketitle

\section{Introduction}
\label{sec:introduction}
For much of today's tree-structured data, treemaps are a versatile and established tool for visualization~\cite{mcnabb2017-surveyofsurveys}.
Since the introduction of treemaps~\cite{johnson1991-spacefilling}, more than 400 publications have been published on this topic.
As for any visualization, it is essential to \enquote{show that [it] is useful in a practical sense}~\cite{kosara2003-userstudies} and \enquote{to understand the potential and limitations}~\cite{plaisant2004-evaluation} of it.
This work examines 69 related user studies from 1992 to 2018 and reviews to what extent these \enquote{provide convincing evidence of utility}~\cite{plaisant2004-evaluation}.
In many of these studies, mistakes in preparation or implementation have been made that led to exaggerated statements. 
As a consequence, \enquote{it is common for claims to be made that go far beyond the results} and \enquote{many of the studies we read are biased}~\cite{ware2013-infovis}.
Using the available data, it is difficult to decide which visualization is best suited for a particular task~\cite{lam2012-empirical}.

With this paper, we contribute a review of these user studies by means of overview tables.
Using the detailed data on the user studies, we conduct an analysis of common pitfalls and discuss possible mitigation.
The full data set is available by means of an interactive community website and public GitHub repository\footnote{\href{https://github.com/varg-dev/treemap-studies}{github.com/varg-dev/treemap-studies}}.
The remainder of this paper is structured as follows.
Section~2 recapitulates the basics of user studies in the field of information visualization.
Section~3 reviews available user studies on treemaps.
Section~4 provides a discussion on the findings and quality of the user studies.
In Section~5, we summarize user study guidelines and publication best-practices.
Section~6 concludes this paper.

\section{Related Work and Methods} %
\label{sec:preliminaries}

Evaluations can examine various aspects of the transformation from data to visual artifacts and the perceptional characteristics of those.
They provide the opportunity to detect shortcomings and pitfalls but also reveal suitable applications and the potential of particular parameterization.
Common evaluation methods are metric-based evaluations, heuristics, or case studies.
User studies are a popular method as they \enquote{offer a scientifically sound method to measure a visualization's performance} and \enquote{can be used to evaluate the strengths and weaknesses of different visualization techniques}~\cite{kosara2003-userstudies}.
Further, user studies are able to provide data on which approach suits a certain objective best and which properties are more likely to achieve desired results.
It investigates \enquote{why a particular technique is effective} and whether it is \enquote{useful in a practical sense}~\cite{kosara2003-userstudies}.

Similar to this paper, other publications reviewed a body of work to assess user studies and other means of evaluation for visualization techniques.
As main distinction, these reviews do not focus on treemap visualization techniques but filter by publication venue or problem domain.
As such, Lam et~al. tagged 361 papers on user studies in information visualization and derived seven scenario categories~\cite{lam2012-empirical}.
These categories were extended by reviewing 581 user studies~\cite{isenberg2013-review}.
Resulting from this, the authors concluded that the research community focuses on \enquote{evaluations of algorithmic performance [\ldots] and qualitative result inspections [\ldots] through images}~\cite{isenberg2013-review} with a lack on \emph{communication through visualization} and \emph{collaborative data analysis} for their initial set of publications.
With a focus on the software visualization domain, Merino et~al. provided a review on software evaluation and visualization techniques covering 181 evaluation papers, whereof 81 papers documented evaluation with user feedback~\cite{merino2018-review}.

\subsection{Methods for User Evaluation}
Users can participate in various ways in an evaluation process and, depending on the type of participation, provide valuable insights into various aspects of a visualization.
Researchers have to choose suitable evaluation techniques depending on their intentions and available means.
The different aspects of evaluation result in multiple independent categorizations of user studies.
Consequently, it is advisable to consciously determine methodology and content for carrying out an evaluation, as proposed by Leroy~et~al.~\cite{leroy2011-designing} and Plaisant~\cite{plaisant2004-evaluation}, respectively. 
Alternative categorizations for information visualization evaluation were presented by Lam et al.~\cite{lam2012-empirical} and its revision by Isenberg et al.~\cite{isenberg2013-review}, which offer a more holistic approach.
As an appropriate analysis would go beyond the scope of this report, this review concentrates on the former two, more focused categorizations.

\paragraph{\textbf{Types of User Studies for Software Evaluation}}
Gondy Leroy suggested a summary of most common evaluation types for software evaluation~\cite{leroy2011-designing}.
Although researchers may use \emph{case studies}, \emph{field studies}, or \emph{descriptive studies} to explain and answer difficult questions, they are not necessarily representative~\cite{plaisant2004-evaluation}. To obtain reproducible results, controlled experiments or their moderated variants are needed:
\par
\emph{Correlation studies} describe changes in certain variables. 
The intention is to identify where a change in one variable coincides with change in other.
They cannot detect the direction of causation though.
\emph{Controlled experiments} evaluate the impact, benefit or drawbacks of information systems.
They are often referred to as demonstration studies, randomized controlled trials, or comparative studies.
Researchers control one or more variables to measure values of one or multiple outcomes. 
This allows to validate causation.
If a random assignment of participants to experiment conditions is not possible---for example due to time, geographical, or social constraints---a \emph{quasi~experiment} is conducted instead.
The chosen type defines the transferability of the user study results to a more general level and whether causation could be shown.

\paragraph{\textbf{Evaluation Practices}} %
Catherine Plaisant~\cite{plaisant2004-evaluation} presented a categorization for the user evaluation of information visualization.
The four categories provide insights on what is compared in which environment (comparison method).
First, \emph{controlled experiments comparing different techniques} are user studies that compare two or more different techniques.
These are referred to as \emph{visualization techniques comparisons} (VTC) and are typically used to contrast new techniques with existing ones.
The second category is \emph{controlled experiment comparing design elements} of an individual visualization technique.
Such studies compare a parameterization with another configuration, e.g., mappings of information to visual display.
For treemaps, we consider three sub-categories: \emph{layout}, \emph{visual variables}, and \emph{configuration} of treemap.
The third category is \emph{usability evaluation}. 
It can provide feedback on a certain visualization tool and prioritize its subsequent improvement and refinement.
Finally and least common, \emph{case studies} enable evaluation in realistic settings, and, similar to Leroy's categorization, facilitate to demonstrate in-context usefulness.
The results of a case study, however, may not be replicable and generalizable~\cite{plaisant2004-evaluation}. Such studies are therefore not included.
During our review we found a number of user studies that do not fit into the categories.
They all deal with treemaps on a meta level in relation to communication and aesthetic perception, or more precisely visual quality, and are assigned to a \emph{meta} category.

\subsection{Design and Conduct of User Studies}%
If researchers want to prove a causation, they have to conduct a controlled experiment. 
Especially for such experiments, possible results can become unusable because of smallest errors.
Mistakes often occur during the preparation or execution of a user study that result in incorrect statements.
Published results of erroneous user studies hardly represent added value for the research community. %
Regarding this, Colin Ware states that \enquote{it is common for claims to be made that go far beyond the results}~\cite{ware2013-infovis}, e.g., generalizing satisfaction statements from single participants \cite{isenberg2013-review}.
Researchers are often unaware of the specifics of those errors and mistakes in the preparation and conduct of user studies. %
Consequently, \enquote{many of the studies we read are biased}~\cite{ware2013-infovis}.
The main objective should therefore be to conduct informative, value-adding and qualitatively meaningful user studies that meet critical validation priorities.

Design problems often emerge from a discrepancy between the research goal and the experiment actually conducted.
Typical, critical design decisions are determining the independent variables and tasks participants have to perform.
Further difficulties are the formulation of hypotheses that fit the research goal, the establishment of appropriate control conditions, and the occurrence of interfering variables~\cite{ware2013-infovis}.
In the broad field of statistical analysis, researchers might lose track or are not aware of implications and exact meaning of results.
Despite methodological sources of error, some user studies lack appropriate and fair control conditions---e.g., another visualization that does not implement a feature in a comparable way---and cause an a-priori biased result.
On top of that, the expectation of the research community has an influence on the way the results of an experiment are published:
it \enquote{is easier to publish results that confirm a hypothesis than results showing no effects}~\cite{ware2013-infovis}.
However, this promotes claims \enquote{that go far beyond the results} and seems to be a common, though certainly unintended, practice.

\subsection{What We Should Expect from User Studies}
The evidence of utility and a deep knowledge of tasks or use cases in relation to different visualization techniques are of major importance.
Well thought-out concepts, careful conduct, and fair control conditions are required.
Researchers should (1) guarantee comprehensibility of their studies by describing their thoughts and procedure precisely, (2) publish their data to support validation and verification of their results, as well as (3) formulate results honestly and distinguish clearly from claims and assumptions.

Study replications are essential to confirm results in a validation process but those are rare~\cite{isenberg2013-review,ware2013-infovis}.
In addition, researchers should evaluate not only their tool, but also independently evaluate the underlying visualization technique and used technology.
For this, however, it is essential to understand any implementation of a visualization as a specialized instance of that visualization (not the visualization itself).
The utility of a specific treemap implementation, for example, might be strongly effected by its aesthetics, responsiveness, colors and shading, size, use of additional graphical effects, embedding in a GUI and the GUI itself, peripherals, and many other aspects.
These, however, are aspects that are not necessarily an inherent part of a treemap visualization technique.
If a sound distinction or exclusion of such aspects is not feasible, the goals of respective studies should be reviewed exhaustively and awareness of respective factors should be increased.
\par
In order to assess the power of treemaps, derive qualified statements from previous results, and judge the research area---regarding topic coverage, clusters, and outliers---a model assessing visualization techniques from multiple usage perspectives is helpful.
Since there are no widely used models on evaluation topics in information visualization yet~\cite{zuk2006-heuristics}, we propose the following set:
\begin{description}
    \item[Use Cases] refer to specific tasks, data domains, data structures, and research questions subject to the visualization technique.
    \item[Perception] refers to human perception and dependencies within visual encodings.
    \item[Interaction] refers to the interaction metaphors applied and their technical implementations and constraints.
    \item[Communication] refers to the language and metaphors used to introduce and talk about visualization techniques and tasks.
    \item[Configuration] refers to specific choice of values for a parameterized visualization technique.
\end{description}
These concepts relates as follows. To solve a \emph{communicated} task within an \emph{use case} scenario, a user \emph{perceives} the visual output of a properly \emph{configured} information visualization and \emph{interacts} with it.

\section{User Studies on Treemaps} %
\label{sec:summary}

\begin{table*}[htbp]
  \raggedright
  \scriptsize%
  \noindent%
  \settowidth{\largestheaderlength}{Controlled Experiment}%
  \renewcommand{\arraystretch}{1.01}
  \rowcolors{3}{white}{gray!15}%
  \setlength{\tabcolsep}{1.9pt}%
  \setlength{\aboverulesep}{0pt}%
  \setlength{\belowrulesep}{0pt}%
  \setlength{\arrayrulewidth}{0.4pt}%
  \caption{Excerpt of treemap survey, part one---including an assessment on study categorization, specific evaluation details, characteristics of the used data set, and chosen treemap configuration (\OK).\newline{}
  Notation: evaluation with (\significOK)\,/\,without (\notsignificOK) statistically significant results reg. treemaps\enspace{}|\enspace{}{\color{gray!40}$\blacksquare$} cell not applicable\enspace{}|\enspace{}{--} unknown\newline{}Abbreviations: Dimension~(\emph{Dim.})\enspace{}|\enspace{}Implementation~(\emph{Impl.})\enspace{}|\enspace{}Topology Depiction~(\emph{Topology})\enspace{}|\enspace{}\newline{}Visualization Technique Comparison~(\emph{VTC})\enspace{}|\enspace{}Visual~(\emph{Vis.})\strut\vspace{-0.75em}}%
  \begin{tabular}{                         r|c!{\tls}c!{\tls}c!{\tls}c!{\tls}c!{\tls}c                                                                                                                                      |c!{\tls}c                                               |c!{\tls}c!{\tls}c                                                      |c!{\tls}c!{\tls}c!{\tls}c!{\tls}c                                                                                         |c!{\tls}c!{\tls}c                                                                  |c!{\tls}c!{\tls}c!{\tls}c                                                                |c!{\tls}c!{\tls}c                                                           |c!{\tls}c!{\tls}c!{\tls}c                                                                                          |c!{\tls}c!{\tls}c                                                            |c!{\tls}c!{\tls}c                                         |c!{\tls}c!{\tls}c!{\tls}c                                                                                                            |}
                                            & \mc{6}{Category}                                                                                                                                                              & \mc{13}{Evaluation}                                                                                                                                                                                                                                                                                                                           & \mc{11}{Dataset(s)}                                                                                                                                                                                                                                                                      & \mc{10}{Treemap Configuration}                                                                                                                                                                                                                                               \\
                                            & \mct{6}{Evaluation Practice}                                                                                                                                                  & \mct{2}{Data}                                          & \mct{3}{Impl.}                                                        & \mct{5}{\#Participants}                                                                                                  & \mct{3}{Measure}                                                                  & \mct{4}{Topology}                                                                       & \mct{3}{Origin}                                                            & \mct{4}{Size (\#n)}                                                                                               & \mct{3}{Shape}                                                              & \mct{3}{Dim.}                                            & \mct{4}{Topology}                                                                                                                   \\
    \rott{\parbox{\largestheaderlength}{~}} & \rott{\strut VTC} & \rott{\strut Design: Layout} & \rott{\strut Design: Vis. Variable} & \rott{\strut Design: Configuration} & \rott{\strut Usability} & \rott{\strut (Meta)} & \rott{\strut quantitative} & \rott{\strut qualitative} & \rott{\strut laboratory} & \rott{\strut online} & \rott{\strut MTurk} & \rott{\strut $<10$} & \rott{\strut $10$--$29$} & \rott{\strut $30$--$49$} & \rott{\strut $50$--$100$} & \rott{\strut $>100$} & \rott{\strut effectiveness} &\rott{\strut efficiency} &\rott{\strut satisfaction} & \rott{\strut list} & \rott{\strut hierarchy} & \rott{\strut tree} & \rott{\strut graph} & \rott{\strut artifical} & \rott{\strut adapted} & \rott{\strut real-world} & \rott{\strut $<10$} & \rott{\strut $10$--$100$} & \rott{\strut $100$--$1000$} & \rott{\strut $>1000$} & \rott{\strut rectangluar} & \rott{\strut circular} & \rott{\strut map-like} & \rott{\strut 2D} & \rott{\strut 2.5D} & \rott{\strut 3D} & \rott{\strut non-nested} & \rott{\strut non-nested: \emph{cushion}} & \rott{\strut nested}  & \rott{\strut nested: \emph{titlebar}} \\\hline
    % table generated from studies.json in treemap-evaluation-data repository
    Directory Browsing~\cite{turo1992-improve,johnson1993-categorical}&\significOK&&&&&&\OK&&\OK&&&&\OK&&&&\significOK&\significOK&\significOK&&&\OK&&&&\OK&&\OK&&&\OK&&&\OK&&&&&\OK&\\
GEnie Experiment~\cite{turo1992-improve}&\significOK&&&\notsignificOK&&&\OK&&\OK&&&&\OK&&&&&\significOK&&&&\OK&&&&\OK&\OK&&&&\OK&&&\OK&&&&&&\OK\\
US Budget Experiment~\cite{johnson1993-categorical}&\significOK&&&&&&\OK&\OK&\OK&&&&&\OK&&&\significOK&\significOK&\notsignificOK&&&\OK&&&&\OK&&\OK&&&\OK&&&\OK&&&&&\OK&\\
3D Information Visualizations~\cite{wiss1999-tasksupport}&\significOK&&&&\significOK&&\OK&&\OK&&&&\OK&&&&\significOK&\significOK&\notsignificOK&&&\OK&&\OK&&&\OK&&&&\OK&&&&&\OK&&&\OK&\\
Sunburst vs. Treemap (small)~\cite{stasko2000-structures}&\significOK&&&&&&\OK&\OK&\OK&&&&&\OK&&&\significOK&\notsignificOK&\notsignificOK&&&\OK&&&\OK&&&\OK&&&\OK&&&\OK&&&\OK&&&\\
Sunburst vs. Treemap (large)~\cite{stasko2000-structures}&\significOK&&&&&&\OK&\OK&\OK&&&&\OK&&&&\notsignificOK&\notsignificOK&\significOK&&&\OK&&&\OK&&&&\OK&&\OK&&&\OK&&&\OK&&&\\
Topology \& Node Size~\cite{barlow2001-comparison}&\significOK&&&&&&\OK&&\OK&&&&\OK&&&&\significOK&\significOK&\significOK&&&\OK&&\OK&&&\OK&\OK&&&\OK&&&\OK&&&&&\OK&\\
Compare to Beamtrees~\cite{vanHam2003-beamtrees}&\significOK&&&&&&\OK&&\OK&&&&\OK&&&&\significOK&\significOK&&&&\OK&&\OK&&&&\OK&\OK&&\OK&&&\OK&&&&\OK&\OK&\\
Compare 5 Tree Visualizations~\cite{kobsa2004-systems}&\significOK&&&\significOK&\significOK&&\OK&\OK&\OK&&&&&\OK&&&\significOK&\significOK&\significOK&&&\OK&&&&\OK&&&\OK&&\OK&&&\OK&&&&\OK&&\OK\\
Treemaps \& Tables~\cite{goldberg2005-monitoring}&\significOK&&&\significOK&&&\OK&\OK&\OK&&&&\OK&&&&&\significOK&\significOK&&&\OK&&&&\OK&\OK&\OK&&&\OK&&&\OK&&&&&\OK&\OK\\
Crop Circles~\cite{wang2006-cropcircles}&\significOK&&&\significOK&\significOK&&\OK&&\OK&&&&\OK&&&&\significOK&\significOK&\notsignificOK&&\OK&&&&&\OK&&&\OK&&\OK&\OK&&\OK&&&&&&\OK\\
Knowledge Discovery~\cite{wang2006-knowledgediscovery}&\significOK&&&&&&\OK&&\OK&&&&\OK&&&&\significOK&\significOK&\significOK&&&\OK&&&&\OK&&&&\OK&\OK&&&\OK&&&&&&\OK\\
Hierarchy Browser~\cite{andrews2007-hvte}&\significOK&&&\notsignificOK&&&\OK&&\OK&&&&&\OK&&&&\significOK&\significOK&&&\OK&&&\OK&&&&\OK&&\OK&&&\OK&&&&&&\OK\\
Search Engine Results~\cite{chu2007-searchengine}&\notsignificOK&&&&&&\OK&\OK&&\OK&&\OK&&&&&&\notsignificOK&\notsignificOK&&&\OK&&&&\OK&&\OK&&&\OK&&&\OK&&&&&&\OK\\
Multiple Views~\cite{teoh2007-multipleviews}&\significOK&&&&\significOK&&&\OK&\OK&&&&\OK&&&&\NA&\NA&\NA&&&\OK&&&&\OK&&\OK&&&\OK&&&\OK&&&&&\OK&\\
Search Result Interface 1~\cite{clarkson2009-resultmaps}&\significOK&&&&&&\OK&&\OK&&&&\OK&&&&\notsignificOK&\notsignificOK&\significOK&&\OK&&&&&\OK&&\OK&&&\OK&&&\OK&&&&&\OK&\\
Search Result Interface 2~\cite{clarkson2009-resultmaps}&\significOK&&&&&&\OK&&\OK&&&&&\OK&&&\notsignificOK&\notsignificOK&\significOK&&\OK&&&&&\OK&&\OK&\OK&&\OK&&&\OK&&&&&\OK&\\
Data Density~\cite{kong2010-guidelines}&\significOK&&&&&&\OK&&&&\OK&&\OK&&&&\significOK&\significOK&&&&\OK&&\OK&&&&\OK&\OK&&\OK&&&\OK&&&\OK&&&\\
Node Link Treemaps~\cite{linsen2011-nodelink}&\notsignificOK&&&&&&\OK&&\OK&&&&\OK&&&&\notsignificOK&\notsignificOK&&&&\OK&&&\OK&&&&\OK&&\OK&&&&\OK&&&\OK&&\\
Software Systems as Cities~\cite{wettel2011-softwarecities}&\OK&&&&&&\OK&&\OK&&&&&\OK&&&\significOK&\significOK&\notsignificOK&&&\OK&&&\OK&&&&\OK&&\OK&&&&\OK&&&&\OK&\\
Eye-tracking for EER-Tutor~\cite{mathews2012-eyetracking}&\significOK&&&&&&\OK&\OK&\OK&&&&\OK&&&&&\significOK&\significOK&\UNKNOWN&\UNKNOWN&\UNKNOWN&\UNKNOWN&&&\OK&\OK&&&&\OK&&&\OK&&&&&&\OK\\
Rapid Overview~\cite{mueller2013-eyetracking}&\significOK&&&&&&\OK&\OK&\OK&&&&&&\OK&&\significOK&\notsignificOK&&&&\OK&&\OK&&&\OK&\OK&&&\OK&&&\OK&&&&&&\OK\\
Software Landscape~\cite{fittkau2015-softwarelandscapes}&\significOK&&&&&&\OK&\OK&\OK&&&&\OK&&&&\notsignificOK&\significOK&\notsignificOK&&&\OK&&&&\OK&&\OK&&&\OK&&&\OK&&&&&\OK&\\
Treemap vs. Euler~\cite{sathiyanarayanan2015-euler}&\notsignificOK&&&&&&\OK&&\OK&&&&\OK&&&&&\notsignificOK&\notsignificOK&&&&\OK&\OK&&&\OK&&&&\OK&&&\OK&&&&&\OK&\\
Treemap vs. Euler 2~\cite{sathiyanarayanan2015-graph}&\notsignificOK&&&&&&\OK&&\OK&&&&\OK&&&&\notsignificOK&\notsignificOK&&&&&\OK&\OK&&&\OK&&&&\OK&&&\OK&&&&&\OK&\\
Photo: Layout~\cite{guerra-gomez2016-photo}&\significOK&&&&&&\OK&&&&\OK&&&&\OK&&\significOK&\significOK&\significOK&\OK&&&&\OK&&&\OK&&&&\OK&&&\OK&&&\OK&&&\\
Compare 3 Tree Visualizations~\cite{mueller2017-perception}&\significOK&&&&\significOK&&\OK&\OK&\OK&&&&&&\OK&&\significOK&\notsignificOK&&&&\OK&&\OK&&&\OK&\OK&&&\OK&&&\OK&&&&&&\OK\\
Wrapped Bars: Comparison~\cite{yalcin2017-wrappedbars}&\significOK&&&&&&\OK&&&&\OK&&\OK&&&&\significOK&\notsignificOK&&\OK&&&&\OK&&&\OK&\OK&&&\OK&&&\OK&&&\OK&&&\\
Wrapped Bars: Ranking~\cite{yalcin2017-wrappedbars}&\significOK&&&&&&\OK&&&&\OK&&&\OK&&&\significOK&\significOK&&\OK&&&&\OK&&&\OK&\OK&&&\OK&&&\OK&&&\OK&&&\\
Wrapped Bars: Distribution~\cite{yalcin2017-wrappedbars}&\significOK&&&&&&\OK&&&&\OK&&\OK&&&&\significOK&\significOK&&\OK&&&&\OK&&&\OK&\OK&&&\OK&&&\OK&&&\OK&&&\\
Pie Charts: Proportion~\cite{huestegge2018-processes}&\significOK&&&&&&\OK&&\OK&&&&\OK&&&&\significOK&\significOK&&\OK&&&&\OK&&&\OK&&&&\OK&&&\OK&&&\OK&&&\\
Pie Charts: Comparison~\cite{huestegge2018-processes}&\significOK&&&&&&\OK&&\OK&&&&\OK&&&&\significOK&\significOK&&\OK&&&&\OK&&&\OK&&&&\OK&&&\OK&&&\OK&&&\\
Pie Charts: Labeling~\cite{huestegge2018-processes}&\significOK&&&&&&\OK&&\OK&&&&\OK&&&&\significOK&\significOK&&\OK&&&&\OK&&&\OK&&&&\OK&&&\OK&&&\OK&&&\\
Pie Charts \& Stacked Bars~\cite{huestegge2018-processes}&\significOK&&&&&&\OK&&\OK&&&&\OK&&&&\significOK&\notsignificOK&\notsignificOK&\OK&&&&\OK&&&\OK&&&&\OK&&&\OK&&&\OK&&&\\\hline
Compare Layout Algorithms~\cite{bederson2002-quantum}&&\significOK&&&&&\OK&&\OK&&&&\OK&&&&&\significOK&\significOK&\OK&&&&\OK&&&&\OK&&&\OK&&&\OK&&&\OK&&&\\
Relative Direction Change~\cite{hahn2017-directionchange}&&\significOK&&&&&\OK&&\OK&&&&\OK&&&&&\significOK&&&&\OK&&&&\OK&\OK&&&&\OK&&&\OK&&&&&\OK&\\
Layout in Cartography~\cite{zhou2017-cartography}&&\significOK&&&&&\OK&&&\OK&&&&&\OK&&\significOK&\significOK&&&&\OK&&&&\OK&\OK&\OK&&&\OK&&&\OK&&&&&\OK&\\\hline
Luminance (Pilot)~\cite{kong2010-guidelines}&&&\significOK&&&&\OK&&&&\OK&&&\OK&&&\significOK&&&&&\OK&&\OK&&&\OK&&&&\OK&&&\OK&&&\OK&&&\\
Aspect Ratio~\cite{kong2010-guidelines}&&&\significOK&&&&\OK&&&&\OK&&\OK&&&&\significOK&&&\NA&\NA&\NA&\NA&\NA&\NA&\NA&\NA&\NA&\NA&\NA&\OK&&&\OK&&&\NA&\NA&\NA&\NA\\
Rectangular Orientation~\cite{liang2012-orientation}&&&\significOK&&&&\OK&&\OK&&&&\OK&&&&\significOK&\notsignificOK&&\NA&\NA&\NA&\NA&\NA&\NA&\NA&\NA&\NA&\NA&\NA&\OK&&&\OK&&&\NA&\NA&\NA&\NA\\
Sketchiness: Perceived Height~\cite{limberger2016-sketchy}&&&\significOK&&&&\OK&&&\OK&&&&&\OK&&\significOK&&&\NA&\NA&\NA&\NA&\NA&\NA&\NA&\NA&\NA&\NA&\NA&\OK&&&&\OK&&\NA&\NA&\NA&\NA\\
Sketchiness: Pre-Attentiveness~\cite{limberger2016-sketchy}&&&\significOK&&&&\OK&&\OK&&&&\OK&&&&\significOK&&&&&\OK&&\OK&&&\OK&\OK&&&\OK&&&&\OK&&&&\OK&\\
Sketchiness: Order~\cite{limberger2016-sketchy}&&&\significOK&&&&\OK&&&\OK&&&&&\OK&&\significOK&&&\NA&\NA&\NA&\NA&\NA&\NA&\NA&\NA&\NA&\NA&\NA&\OK&&&&\OK&&\NA&\NA&\NA&\NA\\
Color Scheme: Model~\cite{xie2017-coloraesthetics}&&&\significOK&&&&\OK&&\OK&&&&\OK&&&&\significOK&&&\NA&\NA&\NA&\NA&\NA&\NA&\NA&\NA&\NA&\NA&\NA&\OK&&&\OK&&&\NA&\NA&\NA&\NA\\
Color Scheme: Aesthetics~\cite{xie2017-coloraesthetics}&&&\significOK&&&&\OK&&\OK&&&&&\OK&&&\significOK&&&\OK&&&&\OK&&&\OK&\OK&&&\OK&&&\OK&&&\OK&&&\\
Color Scheme: Size~\cite{xie2017-coloraesthetics}&&&\significOK&&&&\OK&&\OK&&&&&\OK&&&\significOK&\significOK&&\OK&&&&\OK&&&\OK&&&&\OK&&&\OK&&&\OK&&&\\
Color Scheme: Combination~\cite{xie2017-coloraesthetics}&&&\significOK&&&&\OK&&\OK&&&&&\OK&&&\significOK&&&\OK&&&&\OK&&&\OK&\OK&&&\OK&&&\OK&&&\OK&&&\\\hline
2D and 2.5D~\cite{bladh2004-threedimensions}&&&&\significOK&&&\OK&&\OK&&&&\OK&&&&\significOK&\significOK&&&&\OK&&\UNKNOWN&\UNKNOWN&\UNKNOWN&&&\OK&&\OK&&&\OK&\OK&&&&\OK&\\
Animation in 2.5D~\cite{bladh2005-transitions}&&&&\notsignificOK&&&\OK&\OK&\OK&&&&\OK&&&&\notsignificOK&\notsignificOK&&&&\OK&&\UNKNOWN&\UNKNOWN&\UNKNOWN&&&\OK&&\OK&&&&\OK&&&&\OK&\\
Content Browsing 1~\cite{shi2005-contentbrowsing}&&&&\significOK&&&\OK&&\OK&&&&\OK&&&&&\significOK&&&&\OK&&\OK&&&&\OK&&&\OK&&&\OK&&&\OK&&&\\
Content Browsing 2~\cite{shi2005-contentbrowsing}&&&&\significOK&&&\OK&&\OK&&&&\OK&&&&\notsignificOK&\significOK&&&&\OK&&\OK&&&\OK&&&&\OK&&&\OK&&&\OK&&&\\
Cushion Treemaps~\cite{irani2006-perception}&&&&\significOK&&&\OK&&\OK&&&&\OK&&&&\significOK&\significOK&\notsignificOK&&&\OK&&&\OK&&&&\OK&&\OK&&&\OK&&&\OK&\OK&&\\
Cushion Treemaps (Structure)~\cite{irani2006-perception}&&&&\significOK&&&\OK&&\OK&&&&\OK&&&&\notsignificOK&\significOK&&&&\OK&&\OK&&&&\OK&&&\OK&&&\OK&&&\OK&\OK&&\\
Image Browser~\cite{gomi2008-imagebrowser}&&&&\significOK&&&\OK&&\OK&&&&\OK&&&&&\significOK&\notsignificOK&&&\OK&&\UNKNOWN&\UNKNOWN&\UNKNOWN&&&\OK&&\OK&&&\OK&&&&&\OK&\\
Balloon Focus~\cite{tu2008-multifocus}&&&&\significOK&&&\OK&&\OK&&&&\OK&&&&\significOK&\significOK&\significOK&&&\OK&&&&\OK&&&\OK&&\OK&&&\OK&&&&&&\OK\\
Photo: Count~\cite{guerra-gomez2016-photo}&&&&\significOK&&&\OK&&&&\OK&&&&&\OK&&\significOK&\significOK&\OK&&&&\OK&&&\OK&&&&\OK&&&\OK&&&\OK&&&\\
Photo: Flickr~\cite{guerra-gomez2016-photo}&&&&\significOK&&&\OK&&&\OK&&\OK&&&&&&&\significOK&\OK&&&&&&\OK&\OK&&&&\OK&&&\OK&&&\OK&&&\\
Treemaps vs. Map-like~\cite{biuk-aghai2017-maplike}&&&&\significOK&&&\OK&&\OK&&&&&\OK&&&\significOK&\significOK&&&&\OK&&&&\OK&\OK&&&&\OK&&\OK&\OK&&&\NA&\NA&\NA&\NA\\
Aggregation: Visual Search~\cite{limberger2017-aggregation}&&&&\significOK&&&\OK&&&\OK&&&\OK&&&&\significOK&\notsignificOK&\significOK&&&\OK&&&&\OK&&\OK&\OK&\OK&\OK&&&&\OK&&&&&\OK\\
Aggregation: Readability~\cite{limberger2017-aggregation}&&&&\notsignificOK&&&\OK&&&\OK&&&\OK&&&&\notsignificOK&&&&&\OK&&\OK&&&&\OK&\OK&\OK&\OK&&&&\OK&&&&&\OK\\
Height References~\cite{limberger2018-filtering}&&&&\significOK&&&\OK&&&\OK&&&\OK&&&&\significOK&\notsignificOK&&&&\OK&&&&\OK&&\OK&\OK&&\OK&&&&\OK&&&&\OK&\\\hline
Flexible Hierarchy Interface~\cite{chintalapani2004-flexible}&&&&&\significOK&&&\OK&\OK&&&\OK&&&&&\NA&\NA&\NA&&\OK&&&&&\OK&\OK&&&&\OK&&&\OK&&&&&&\OK\\
Virtual Reality~\cite{fittkau2015-virtualreality}&&&&&\significOK&&\OK&\OK&\OK&&&&\OK&&&&\NA&\NA&\NA&&&\OK&&&&\OK&\OK&&&&\OK&&&&\OK&&&&&\OK\\
Code Maps 1~\cite{hawes2015-codesurveyor}&&&&&\significOK&&&\OK&\OK&&&\OK&&&&&\NA&\NA&\NA&&&\OK&&&&\OK&&&\OK&\OK&&&\OK&\OK&&&\OK&&&\\
Code Maps 2~\cite{hawes2015-codesurveyor}&&&&&\significOK&&&\OK&\OK&&&\OK&&&&&\NA&\NA&\NA&&&\OK&&&&\OK&&&&\OK&&&\OK&\OK&&&\OK&&&\\\hline
Aesthetic~\cite{cawthon2007-aesthetic}&&&&&&\significOK&\OK&&&\OK&&&&&&\OK&\notsignificOK&\significOK&\significOK&&&\OK&&&&\OK&\OK&&&&\OK&&&\OK&\OK&&&&\OK&\OK\\
Visual Metaphors~\cite{ziemkiewicz2008-metaphors}&&&&&&\significOK&\OK&&\OK&&&&&\OK&&&\significOK&\significOK&&&&\OK&&\OK&&&&\OK&&&\OK&&&\OK&&&&&&\OK\\
Individual Visual Metaphors~\cite{ziemkiewicz2009-preconceptions}&&&&&&\significOK&\OK&&&&\OK&&&&\OK&&\significOK&\significOK&\significOK&&&\OK&&\OK&&&\OK&&&&\OK&&&\OK&&&&&&\OK\\
LoC \& Visual Metaphors~\cite{ziemkiewicz2011-locus}&&&&&&\significOK&\OK&&&&\OK&&&&&\OK&\significOK&\significOK&\significOK&&&\OK&&&&\OK&\OK&&&&\OK&&&\OK&&&&&&\OK\\\hline

  \end{tabular}%
  \label{tab:overview-basic}%
\end{table*}%


\begin{table*}[htbp]
  \raggedright
  \scriptsize%
  \noindent% 
  \settowidth{\largestheaderlength}{Controlled Experiment}%
  \renewcommand{\arraystretch}{1.01}
  \rowcolors{3}{white}{gray!20}%
  \setlength{\tabcolsep}{1.9pt}%
  \setlength{\aboverulesep}{0pt}
  \setlength{\belowrulesep}{0pt}
  \setlength{\arrayrulewidth}{0.4pt}
  \caption{Excerpt of treemap survey, part two---including an assessment of details on chosen treemap configuration, interaction, user interface, and study tasks (\OK).\newline{}
  Notation: {\color{gray!40}$\blacksquare$} table cell not applicable\enspace{}|\enspace{}{--} unknown\enspace{}|\enspace{}\newline{}
  Abbreviations: Verbal Metaphor (\emph{Mtph.})\enspace{}|\enspace{}Dimension (\emph{Dim.})\enspace{}|\enspace{}{}size~(\emph{S})\enspace{}|\enspace{}height~(\emph{H})\enspace{}|\enspace{}color~(\emph{C})\enspace{}|\enspace{}structure~(\emph{R})\enspace{}|\newline{}(sub)tree-attribute~(\emph{T})\enspace{}|\enspace{}count~(\emph{U})\enspace{}|\enspace{}name\,/\,id~(\emph{N})\enspace{}|\enspace{}locate~(\emph{L})\enspace{}|\enspace{}other attribute\,/\,visual variable~(\emph{\rhombusOK }).\strut\vspace{-0.75em}}
  \begin{tabular}{                         r|c!{\tls}c!{\tls}c!{\tls}c                                                              |c!{\tls}c!{\tls}c!{\tls}c                                                                  |c!{\tls}c!{\tls}c!{\tls}c                                                                        |c!{\tls}c!{\tls}c!{\tls}c!{\tls}c!{\tls}c!{\tls}c                                                                                                                                                                |c!{\tls}c!{\tls}c!{\tls}c!{\tls}c!{\tls}c!{\tls}c!{\tls}c!{\tls}c!{\tls}c                                                                                                                                                                                                                    |c!{\tls}c                                         |c!{\tls}c!{\tls}c                                                         |}      
                                            & \mc{19}{Treemap Configuration Details}                                                                                                                                                                                                                                                                                                                                                                                                                                                                & \mc{15}{Tasks}                                                                                                                                                                                                                                                                                                                                                                                                            \\      
                                            & \mct{4}{Offset}                                                                       & \mct{4}{Labeling}                                                                         & \mct{4}{Navigation}                                                                             & \mct{7}{UI Features}                                                                                                                                                                                            & \mct{10}{Types}                                                                                                                                                                                                                                                                             & \mct{2}{Mtph.}                                   & \mct{3}{Dim.}                                                            \\      
    \rott{\parbox{\largestheaderlength}{~}} & \rott{\strut none} & \rott{\strut contour} & \rott{\strut thin} & \rott{\strut broad} & \rott{\strut none} & \rott{\strut categories} & \rott{\strut leaves} & \rott{\strut full} & \rott{\strut none} & \rott{\strut node zooming} & \rott{\strut restricted} & \rott{\strut free} & \rott{\strut none} & \rott{\strut search} & \rott{\strut filtering} & \rott{\strut attribute mapping} & \rott{\strut mapping parameters} & \rott{\strut adjustable offset} & \rott{\strut get attribute values} & \rott{\strut identify value} & \rott{\strut access information} & \rott{\strut sort} & \rott{\strut compare} & \rott{\strut identify distribution} & \rott{\strut distinguish} & \rott{\strut associate} & \rott{\strut correlate} & \rott{\strut create information} & \rott{\strut other} & \rott{\strut containment} & \rott{\strut levels} & \rott{\strut 1D tasks} & \rott{\strut 2D tasks} & \rott{\strut 3D tasks} \\\hline
    % table generated from studies.json in treemap-evaluation-data repository
    Directory Browsing~\cite{turo1992-improve,johnson1993-categorical}&&&\OK&&&&&\OK&\OK&&&&&&&&\OK&\OK&\OK&NSU\rhombusOK&LCR&&&&S&&&&&\OK&&\OK&\OK&\\
GEnie Experiment~\cite{turo1992-improve}&&\OK&&&&&&\OK&\OK&&&&&&&&&&\OK&\UNKNOWN&\UNKNOWN&\UNKNOWN&\UNKNOWN&\UNKNOWN&\UNKNOWN&\UNKNOWN&\UNKNOWN&\UNKNOWN&\UNKNOWN&\UNKNOWN&\UNKNOWN&\UNKNOWN&\UNKNOWN&\UNKNOWN\\
US Budget Experiment~\cite{johnson1993-categorical}&&&&\OK&&&&\OK&&\OK&&&&&&&\OK&\OK&\OK&NSU&LSR&S&S&&S&&&&&&\OK&\OK&\OK&\\
3D Information Visualizations~\cite{wiss1999-tasksupport}&&&&\OK&&&&\OK&&&&\OK&\OK&&&&&&&U&L&R&&&&&&&&\OK&&\OK&&\\
Sunburst vs. Treemap (small)~\cite{stasko2000-structures}&&\OK&&&&&\OK&&&\OK&&&&&&&&&\OK&&LSC&SR&&&SR&\rhombusOK&&&&\OK&&\OK&\OK&\OK\\
Sunburst vs. Treemap (large)~\cite{stasko2000-structures}&&\OK&&&&&\OK&&&\OK&&&&&&&&&\OK&&LSC&SR&&&SR&\rhombusOK&&&&\OK&&\OK&\OK&\OK\\
Topology \& Node Size~\cite{barlow2001-comparison}&&&&\OK&&&&\OK&\OK&&&&\OK&&&&&&&T&LR&S&&&SR&&&&&&\OK&\OK&&\\
Compare to Beamtrees~\cite{vanHam2003-beamtrees}&\OK&&\OK&&\OK&&&&&\OK&&&\UNKNOWN&\UNKNOWN&\UNKNOWN&\UNKNOWN&\UNKNOWN&\UNKNOWN&\UNKNOWN&NT\rhombusOK&LR&S&&&SR&&&&\rhombusOK&&\OK&\OK&&\\
Compare 5 Tree Visualizations~\cite{kobsa2004-systems}&\OK&&\OK&&\OK&&&\OK&&\OK&&&&&\OK&\OK&\OK&\OK&\OK&NUT&LCR\rhombusOK&T&&&SR\rhombusOK&&&&&\OK&\OK&\OK&\OK&\\
Treemaps \& Tables~\cite{goldberg2005-monitoring}&&&\OK&&&&&\OK&\OK&&&&&&&\OK&&&\OK&U\rhombusOK&LS&R&&&S\rhombusOK&&\rhombusOK&\rhombusOK&&\OK&&\OK&\OK&\\
Crop Circles~\cite{wang2006-cropcircles}&&&\OK&&\OK&&&\OK&&\OK&&\OK&&\OK&&&&&\OK&&LR&R&&&R&&&&\rhombusOK&&\OK&\OK&&\\
Knowledge Discovery~\cite{wang2006-knowledgediscovery}&&&\OK&&&&&\OK&&\OK&&&&&\OK&\OK&\OK&\OK&\OK&N&&&&&&\rhombusOK&&&&\NA&\NA&&\OK&\OK\\
Hierarchy Browser~\cite{andrews2007-hvte}&&&\OK&&&&&\OK&&\OK&&&\OK&&&&&&&NU&LR&R&&&R&&&&&\OK&&\OK&&\\
Search Engine Results~\cite{chu2007-searchengine}&&&&\OK&&&&\OK&\OK&&&&&&&\OK&\OK&&&\UNKNOWN&\UNKNOWN&\UNKNOWN&\UNKNOWN&\UNKNOWN&\UNKNOWN&\UNKNOWN&\UNKNOWN&\UNKNOWN&\UNKNOWN&\NA&\NA&\NA&\NA&\NA\\
Multiple Views~\cite{teoh2007-multipleviews}&&&&\OK&&&&\OK&\OK&&&&\OK&&&&&&&&&&&&&&&\rhombusOK&&\NA&\NA&\OK&\OK&\\
Search Result Interface 1~\cite{clarkson2009-resultmaps}&&&\OK&&&\OK&&&\OK&&&&\OK&&&&&&&\UNKNOWN&\UNKNOWN&\UNKNOWN&\UNKNOWN&\UNKNOWN&\UNKNOWN&\UNKNOWN&\UNKNOWN&\UNKNOWN&\UNKNOWN&\NA&\NA&\UNKNOWN&\UNKNOWN&\UNKNOWN\\
Search Result Interface 2~\cite{clarkson2009-resultmaps}&&&\OK&&&\OK&&&\OK&&&&\OK&&&&&&&\UNKNOWN&\UNKNOWN&\UNKNOWN&\UNKNOWN&\UNKNOWN&\UNKNOWN&\UNKNOWN&\UNKNOWN&\UNKNOWN&\UNKNOWN&\NA&\NA&\UNKNOWN&\UNKNOWN&\UNKNOWN\\
Data Density~\cite{kong2010-guidelines}&&\OK&&&\OK&&&&\OK&&&&\OK&&&&&&&&&S&S&&&&&&&\NA&\NA&\OK&&\\
Node Link Treemaps~\cite{linsen2011-nodelink}&\OK&&&&\OK&&&&&&&\OK&&&&&&&\OK&NUT&LR\rhombusOK&T&&&SR&&&&&\OK&\OK&\OK&\OK&\\
Software Systems as Cities~\cite{wettel2011-softwarecities}&&&&\OK&\OK&&&&&&\OK&&&&\OK&\OK&&\OK&\OK&&C\rhombusOK&\rhombusOK&&&HR\rhombusOK&&&\rhombusOK&\rhombusOK&\NA&\NA&\OK&\OK&\OK\\
Eye-tracking for EER-Tutor~\cite{mathews2012-eyetracking}&&\OK&&&&&&\OK&\OK&&&&\OK&&&&&&&\UNKNOWN&\UNKNOWN&\UNKNOWN&\UNKNOWN&\UNKNOWN&\UNKNOWN&\UNKNOWN&\UNKNOWN&\UNKNOWN&\UNKNOWN&\UNKNOWN&\UNKNOWN&\UNKNOWN&\UNKNOWN&\UNKNOWN\\
Rapid Overview~\cite{mueller2013-eyetracking}&&&\OK&&&&&\OK&\OK&&&&\OK&&&&&&&U&R&S&&&&&&\rhombusOK&&&\OK&\OK&&\\
Software Landscape~\cite{fittkau2015-softwarelandscapes}&&&&\OK&&&&\OK&\OK&&&&\OK&&&&&&&N&&&&&&&&&\rhombusOK&\NA&\NA&\NA&\NA&\NA\\
Treemap vs. Euler~\cite{sathiyanarayanan2015-euler}&&&\OK&&&&\OK&&\OK&&&&\OK&&&&&&&\UNKNOWN&\UNKNOWN&\UNKNOWN&\UNKNOWN&\UNKNOWN&\UNKNOWN&\UNKNOWN&\UNKNOWN&\UNKNOWN&\UNKNOWN&\NA&\NA&\OK&&\\
Treemap vs. Euler 2~\cite{sathiyanarayanan2015-graph}&&&\OK&&&&\OK&&\OK&&&&\OK&&&&&&&\UNKNOWN&\UNKNOWN&\UNKNOWN&\UNKNOWN&\UNKNOWN&\UNKNOWN&\UNKNOWN&\UNKNOWN&\UNKNOWN&\UNKNOWN&\NA&\NA&\OK&&\\
Photo: Layout~\cite{guerra-gomez2016-photo}&&\OK&&&\OK&&&&\OK&&&&\OK&&&&&&&&L&&&&S&&&&&\NA&\NA&\NA&\NA&\NA\\
Compare 3 Tree Visualizations~\cite{mueller2017-perception}&&&\OK&&&&&\OK&\OK&&&&\OK&&&&&&&U&R&S&&&&&&\rhombusOK&&&\OK&\OK&&\\
Wrapped Bars: Comparison~\cite{yalcin2017-wrappedbars}&&\OK&&&\OK&&&&\OK&&&&\OK&&&&&&&&&S&S&&&&&&&\NA&\NA&\OK&&\\
Wrapped Bars: Ranking~\cite{yalcin2017-wrappedbars}&&\OK&&&\OK&&&&\OK&&&&\OK&&&&&&&&&S&&&&&&&&\NA&\NA&\OK&&\\
Wrapped Bars: Distribution~\cite{yalcin2017-wrappedbars}&&\OK&&&\OK&&&&\OK&&&&\OK&&&&&&&&&&&S&&&&&&\NA&\NA&\OK&&\\
Pie Charts: Proportion~\cite{huestegge2018-processes}&&\OK&&&&&\OK&&\OK&&&&\OK&&&&&&&S&&&&&&&&&&\NA&\NA&\OK&&\\
Pie Charts: Comparison~\cite{huestegge2018-processes}&&\OK&&&&&\OK&&\OK&&&&\OK&&&&&&&&&S&&&&&&&&\NA&\NA&\OK&&\\
Pie Charts: Labeling~\cite{huestegge2018-processes}&&\OK&&&&&\OK&&\OK&&&&\OK&&&&&&&&&S&&&&&&&&\NA&\NA&\OK&&\\
Pie Charts \& Stacked Bars~\cite{huestegge2018-processes}&&\OK&&&&&\OK&&\OK&&&&\OK&&&&&&&&&S&&&&&&&&\NA&\NA&\OK&&\\\hline
Compare Layout Algorithms~\cite{bederson2002-quantum}&&\OK&&&&&\OK&&\OK&&&&\OK&&&&&&&&L&&&&&&&&&\NA&\NA&\OK&&\\
Relative Direction Change~\cite{hahn2017-directionchange}&&&\OK&&&&\OK&&\OK&&&&\OK&&&&&&&&L&&&&&&&&&\NA&\NA&\NA&\NA&\NA\\
Layout in Cartography~\cite{zhou2017-cartography}&&&\OK&&&&&\OK&\OK&&&&\OK&&&&&&&N&LR&S&&SC&&&\rhombusOK&\rhombusOK&&\OK&\OK&\OK&\OK&\\\hline
Luminance (Pilot)~\cite{kong2010-guidelines}&&\OK&&&\OK&&&&\OK&&&&\OK&&&&&&&&&S&S&&&&&&&\NA&\NA&\OK&&\\
Aspect Ratio~\cite{kong2010-guidelines}&\NA&\NA&\NA&\NA&\NA&\NA&\NA&\NA&\NA&\NA&\NA&\NA&\NA&\NA&\NA&\NA&\NA&\NA&\NA&&&S&S&&&&&&&\NA&\NA&\NA&\NA&\NA\\
Rectangular Orientation~\cite{liang2012-orientation}&\NA&\NA&\NA&\NA&\NA&\NA&\NA&\NA&\NA&\NA&\NA&\NA&\NA&\NA&\NA&\NA&\NA&\NA&\NA&&&S&S&&&&&&&\NA&\NA&\NA&\NA&\NA\\
Sketchiness: Perceived Height~\cite{limberger2016-sketchy}&\NA&\NA&\NA&\NA&\NA&\NA&\NA&\NA&\OK&&&&\NA&\NA&\NA&\NA&\NA&\NA&\NA&&&&H&&&&&&&\NA&\NA&\OK&&\\
Sketchiness: Pre-Attentiveness~\cite{limberger2016-sketchy}&&&\OK&&\OK&&&&\OK&&&&\OK&&&&&&&&&&&&\rhombusOK&&&&&\NA&\NA&\OK&&\\
Sketchiness: Order~\cite{limberger2016-sketchy}&\NA&\NA&\NA&\NA&\NA&\NA&\NA&\NA&\NA&\NA&\NA&\NA&\NA&\NA&\NA&\NA&\NA&\NA&\NA&&&\rhombusOK&&&&&&&&\NA&\NA&\OK&&\\
Color Scheme: Model~\cite{xie2017-coloraesthetics}&\NA&\NA&\NA&\NA&\NA&\NA&\NA&\NA&\NA&\NA&\NA&\NA&\NA&\NA&\NA&\NA&\NA&\NA&\NA&&&S&&&&&&&&\NA&\NA&\OK&&\\
Color Scheme: Aesthetics~\cite{xie2017-coloraesthetics}&\OK&&&&\OK&&&&\OK&&&&\OK&&&&&&&&&&&&&&&&\rhombusOK&\NA&\NA&\NA&\NA&\NA\\
Color Scheme: Size~\cite{xie2017-coloraesthetics}&\OK&&&&&&\OK&&\OK&&&&\OK&&&&&&&&&S&&&&&&&&\NA&\NA&\OK&&\\
Color Scheme: Combination~\cite{xie2017-coloraesthetics}&\OK&&&&\OK&&&&\OK&&&&\OK&&&&&&&&&S&&&&&&&\rhombusOK&\NA&\NA&\OK&&\\\hline
2D and 2.5D~\cite{bladh2004-threedimensions}&&&\OK&&&\OK&&&&&\OK&&&&\OK&&&&\OK&N&LCR&ST&&C&SR&&&&&\OK&&\OK&\OK&\\
Animation in 2.5D~\cite{bladh2005-transitions}&&&\OK&&&\OK&&&&\OK&\OK&&&&\OK&&\OK&&\OK&&LCR&&&C&S&&&&\rhombusOK&\OK&&\OK&\OK&\\
Content Browsing 1~\cite{shi2005-contentbrowsing}&&\OK&&&\OK&&&&&\OK&&&\OK&&&&&&&&L&&&&&&&&&\NA&\NA&\OK&&\\
Content Browsing 2~\cite{shi2005-contentbrowsing}&&\OK&&&\OK&&&&&\OK&&&\OK&&&&&&&&L&&&&&&&&&\UNKNOWN&\UNKNOWN&\OK&&\\
Cushion Treemaps~\cite{irani2006-perception}&\OK&\OK&&&\OK&&&&&\OK&&&\OK&&&&&&&NU&LCR&&&C&SR&&&&&\OK&&\OK&\OK&\\
Cushion Treemaps (Structure)~\cite{irani2006-perception}&\OK&\OK&&&\OK&&&&\OK&&&&\OK&&&&&&&&&&&&&&&&\rhombusOK&\UNKNOWN&\UNKNOWN&\OK&&\\
Image Browser~\cite{gomi2008-imagebrowser}&&&\OK&&\OK&&&&&&&\OK&\OK&&&&&&&&L&&&&&&&&&\UNKNOWN&\UNKNOWN&\OK&&\\
Balloon Focus~\cite{tu2008-multifocus}&&&\OK&&&\OK&&&\OK&&&&&&&&\OK&&\OK&U&C\rhombusOK&&&&S&&&&\rhombusOK&\NA&\NA&\OK&&\\
Photo: Count~\cite{guerra-gomez2016-photo}&&\OK&&&\OK&&&&\OK&&&&\OK&&&&&&&&L&&&&&&&&&\NA&\NA&\NA&\NA&\NA\\
Photo: Flickr~\cite{guerra-gomez2016-photo}&&\OK&&&\OK&&\OK&&\OK&&&&\OK&&&&&&&&&&&&&&&&\rhombusOK&\UNKNOWN&\UNKNOWN&\UNKNOWN&\UNKNOWN&\UNKNOWN\\
Treemaps vs. Map-like~\cite{biuk-aghai2017-maplike}&&\OK&&&&&&\OK&\OK&&&&\OK&&&&&&&&&&&&&&\rhombusOK&&&\NA&\NA&\OK&&\\
Aggregation: Visual Search~\cite{limberger2017-aggregation}&&&\OK&&&\OK&&&\OK&&&&\OK&&&&&&&&&&&&S\rhombusOK&&&&&\NA&\NA&\OK&\OK&\\
Aggregation: Readability~\cite{limberger2017-aggregation}&&&\OK&&&\OK&&&\OK&&&&\OK&&&&&&&\rhombusOK&&&&&&&\rhombusOK&&&\NA&\NA&\OK&\OK&\\
Height References~\cite{limberger2018-filtering}&&&\OK&&\OK&&&&\OK&&&&&&\OK&&&&&&&&&&H&\rhombusOK&&&&\NA&\NA&\OK&&\\\hline
Flexible Hierarchy Interface~\cite{chintalapani2004-flexible}&&&\OK&&&&&\OK&\OK&&&&&&\OK&\OK&\OK&\OK&\OK&N\rhombusOK&\rhombusOK&&S&&S&\rhombusOK&&&&\NA&\NA&\OK&\OK&\\
Virtual Reality~\cite{fittkau2015-virtualreality}&&&\OK&&&&&\OK&&&&\OK&\OK&&&&&&&N&LR&&&&S&&&\rhombusOK&\rhombusOK&\OK&&\OK&&\\
Code Maps 1~\cite{hawes2015-codesurveyor}&&\OK&&&&\OK&&&\OK&&&&\OK&&&&&&&&&&&&&&\rhombusOK&&\rhombusOK&\NA&\NA&\NA&\NA&\NA\\
Code Maps 2~\cite{hawes2015-codesurveyor}&&\OK&&&&&&\OK&&&&\OK&&\OK&&&&&&&L&&&&&&&&\rhombusOK&\NA&\NA&\NA&\NA&\NA\\\hline
Aesthetic~\cite{cawthon2007-aesthetic}&&\OK&\OK&&\UNKNOWN&\UNKNOWN&\UNKNOWN&\UNKNOWN&\OK&&&&\OK&&&&&&&\UNKNOWN&\UNKNOWN&\UNKNOWN&\UNKNOWN&\UNKNOWN&\UNKNOWN&\UNKNOWN&\UNKNOWN&\UNKNOWN&\UNKNOWN&\UNKNOWN&\UNKNOWN&\UNKNOWN&\UNKNOWN&\UNKNOWN\\
Visual Metaphors~\cite{ziemkiewicz2008-metaphors}&&&\OK&&&&&\OK&\OK&&&&\OK&&&&&&&NU&LCR&&&&R&&&&&\OK&\OK&\OK&\OK&\\
Individual Visual Metaphors~\cite{ziemkiewicz2009-preconceptions}&&&\OK&&&&&\OK&\OK&&&&\OK&&&&&&&&LR&RT&&&&&&&&\OK&\OK&\OK&&\\
LoC \& Visual Metaphors~\cite{ziemkiewicz2011-locus}&&&&\OK&&&&\OK&\OK&&&&\OK&&&&&&&&LS\rhombusOK&&&&&\rhombusOK&&&\rhombusOK&\OK&\OK&\OK&&\\\hline

  \end{tabular}%
  \label{tab:overview-details}%
\end{table*}%


The evaluation of treemaps started almost simultaneously with their invention in 1991~\cite{turo1992-improve,johnson1993-categorical}.
In contrast to the initial definition of treemaps as a recursive subdivision of a 2D surface~\cite{johnson1991-spacefilling}, hereinafter treemaps denote visualizations of data based on parent-child-relations where child nodes are located within their parent surfaces but are not required to utilize all of the available parent node's space~\cite{bederson2002-quantum}.
Thereby, treemaps include not only subdivisions of 2D rectangles but the extension to $n$-dimensional spaces~\cite{johnson1993-categorical} and other shapes~\cite{wattenberg2005-spacefillingcurves,balzer2005-voronoi,auber2013-gospermap}.
Since the initial publication in 1991, over 400 publications regarding treemaps were published in conference proceedings and journal articles of respective publishers, i.e., IEEE (TVCG, iV, VISSOFT, VIS, INFOVIS, and VAST), ACM (ToG, VINCI, OzCHI, ICSE, and AVI), Springer (CompStat, ISWC, ITS, and APCHI), EG (CGF, EuroVis, and EuroRVVV), SciTePress (VISIGRAPP\,/\,IVAPP), and SagePub (IVI).
For the list of papers considered in this review, we extended this corpus by (1) selected doctoral theses and technical reports as some user studies were not published with peer review, (2) a short list of individual papers by lesser frequented publishers, their journals, and conferences.
We selected papers that include evaluations and studies with participants, whereby the study is classifiable as \emph{Summative User Study (Evaluation)}~\cite{munzner2008-pitfalls}.
If multiple user studies are documented within one publication, they are treated as different studies within the review and analysis.
We identified 69 user studies in according publications on treemaps until 2018 and considered all of them for structural analysis and 63 of them for the evaluation of the results.
Thereby, we excluded case studies as they are hard to reproduce and the statement generalization is limited~\cite{plaisant2004-evaluation}.
More details on the excluded studies is provided in Section~\ref{sec:summary:findings}. 
The per-year distribution is shown in Figure~\ref{fig:studycount}.
\begin{figure}[t]%
	\centering%
	

\definecolor{plots-blue}{RGB}{77,115,167}%
\makeatletter
\newcommand\printtwodigits[1]{\two@digits{#1}}
\makeatother
\begin{tikzpicture}[
]
\begin{axis}[name=plot-main,
	ybar,
	xmin=1992,
	xmax=2018,
	ymin=0,
	ymax=13,
	axis x line*=bottom,
	hide y axis,
	width=1.17\columnwidth,
	height=2.4cm,
	enlarge x limits=0.04,
	nodes near coords,
	bar width=6pt,
	every node near coord/.style={
	    color=black,
	    font=\scriptsize,
	},
	xtick={1992, ..., 2018},
	xticklabel={
    	\pgfmathparse{mod(\tick,100)}
        \pgfmathtruncatemacro{\shortenedyear}{\pgfmathresult}
        \pgfmathsetbasenumberlength{2}
	    '$\printtwodigits{\shortenedyear}$
    },
	x tick label style={rotate=90,anchor=east,font=\footnotesize},
	/pgf/number format/1000 sep={},
	]
	\addplot+[fill=black,
	draw=none,
	] table[x index=1,y index=0,col sep=semicolon] {data/studies-per-year-summarized.txt};
\end{axis}

\end{tikzpicture}%
%
	\vspace{-1.0em}
	\caption{Number of user studies considered per year. The publication year 2018 is not completely considered.}%
	\label{fig:studycount}%
\end{figure}%

\subsection{Overview of User Studies} %
\label{sec:summary:overview}
The results of the survey are provided as a comprehensive data set of which an excerpt is visualized in Table~\ref{tab:overview-basic} and Table~\ref{tab:overview-details}.
Table~\ref{tab:vtc} contains information on compared visualization techniques.
The complete data set, its schema, and a description for every data category is interactively available on the community website\footnotemark[1] on which researchers are invited to collaborate as well as refine and extend the data set.

In the following, the user studies are described statistically.
Unless otherwise stated, all figures refer to all 69 user studies analyzed:
of the user studies analyzed, almost half contain comparisons with other visualization techniques (\emph{VTC}) and more than one quarter \emph{configuration} aspects.
There are only three studies targeting \emph{layout}, ten studies evaluating \emph{visual variables}, and four studies in the \emph{meta} category.
In \emph{VTC} studies, typically multiple visualizations are compared.
Comparisons with node-link diagrams were the most frequent, followed by sunburst views and lists.
Of the six layout algorithms, only the \emph{Strip} and \emph{Squarified} algorithms appeared in all layout user studies.
The evaluation data type was purely \emph{qualitative} in four studies and purely \emph{quantitative} in more than three quarters of the studies.
The latter use more than 80\,\% of \emph{effectiveness} and \emph{efficiency} measures~\cite{leroy2011-designing} for their dependent variables.
Less than half used \emph{satisfaction} measures.
The size of used data sets ranged broadly from five to 200\,000 nodes with \emph{VTC} studies using the most extensive data sets in average and \emph{meta} studies the smallest.
For studies targeting \emph{visual variables}, most did not include an actual real-world data set.
Almost three quarters of the studies were conducted in a \emph{laboratory} setting, while \emph{online} and \emph{Mechanical Turk} studies accounted for only one eighth each.
The number of participants ranged from five to 285 (Figure~\ref{fig:participants}). %

The most common configuration of treemaps were \emph{2D} \emph{rectangular} treemaps.
More than 80\,\% of the studies included rectangular treemaps, while \emph{circular} and \emph{map-like} shapes were uncommon, with one and three studies, respectively. 
There were only one study covering \emph{3D} treemaps and 12 studies containing \emph{2.5D} treemaps.
More than half of the studies included \emph{non-nested} treemaps~\cite{johnson1993-categorical} (including \emph{cushion} treemaps~\cite{vanwijk1999-cushion}) and almost 40\,\% of the studies \emph{nested} treemaps (including \emph{titlebar} treemaps~\cite{goldberg2005-monitoring}) for topology depiction.
Targeted tasks were categorized according to Cantu et al.~\cite{cantu2017-relations} and the corresponding visual variables or attributes.
The most common task types were \emph{access information} tasks included in more than half of the studies, followed by \emph{sort}, \emph{distinguish}, and \emph{compare} tasks.
Least common task types were \emph{associate}, \emph{correlate}, and \emph{identify distribution} tasks. 
Size-related tasks---part of more than half of the studies---represent the majority of used attributes, followed by tasks related to structure and color.
With 45\,\%, the most common tasks were locate-a-node tasks and sort-by-size tasks with more than one third of the studies.

\begin{figure*}[t]%
	\centering%
	

\pgfdeclareplotmark{mysquare}{%
	\path[draw=black] (-0.5\pgfplotmarksize,-0.5\pgfplotmarksize) rectangle (0.5\pgfplotmarksize,0.5\pgfplotmarksize);%
}%
\pgfdeclareplotmark{mycross}{%
	\path[draw=black] (0.0,0.0) circle (0.5\pgfplotmarksize);%
	\path[draw=black] (-0.5\pgfplotmarksize,-0.5\pgfplotmarksize) -- (0.5\pgfplotmarksize,0.5\pgfplotmarksize);%
	\path[draw=black] (-0.5\pgfplotmarksize,0.5\pgfplotmarksize) -- (0.5\pgfplotmarksize,-0.5\pgfplotmarksize);%
}%
\definecolor{plots-laboratory}{RGB}{106, 169,108}%
\definecolor{plots-online}{RGB}{77,115,167}%
\definecolor{plots-mturk}{RGB}{75,187,187}%
\newlength{\axislabellength}%
\settowidth{\axislabellength}{\small 1}%
\newcommand*{\printmedian}{}
\newcommand*{\participantsplots}{
\addplot+ [boxplot={draw position=14},
my boxplot,
laboratory style,
] table[x index=0,y index=0] {data/participants-laboratory-vtc.csv} \printmedian;
\addplot+[online style, only marks] table[x index=0,y expr=13] {data/participants-online-vtc.csv};
\addplot+ [boxplot={draw position=12},
my boxplot,
mturk style,
] table[x index=0,y index=0] {data/participants-mturk-vtc.csv} \printmedian;
\addplot+ [boxplot={draw position=10},
my boxplot,
laboratory style,
] table[x index=0,y index=0] {data/participants-laboratory-configuration.csv} \printmedian;
\addplot+ [boxplot={draw position=9},
my boxplot,
online style,
] table[x index=0,y index=0] {data/participants-online-configuration.csv} \printmedian;
\addplot+[mturk style, only marks] table[x index=0,y expr=8] {data/participants-mturk-configuration.csv};
\addplot+ [boxplot={draw position=6},
my boxplot,
laboratory style,
] table[x index=0,y index=0] {data/participants-laboratory-usability.csv} \printmedian;
\addplot+[laboratory style, only marks] table[x index=0,y expr=2] {data/participants-laboratory-meta.csv};
\addplot+[online style, only marks] table[x index=0,y expr=1] {data/participants-online-meta.csv};
\addplot+[mturk style, only marks] table[x index=0,y expr=0] {data/participants-mturk-meta.csv};
\newcommand{\xValue}{0}%
\foreach \xValue in {10, 20, 30, 40, 50, 60, 70, 80, 90, 100, 150, 200, 250, 300,350} {%
	\edef\temp{\noexpand\draw[thin,dotted,draw=gray] (axis cs:\xValue,0.5) -- (axis cs:\xValue,14.5);}%
	\temp%
}%
\newcommand{\yValue}{0}%
\foreach \yValue in {3, 7, 11} {%
	\edef\temp{\noexpand\draw[thin,draw=gray] (axis cs:0,\yValue) -- (axis cs:420,\yValue);}%
	\temp
}%
}%
\normalsize
\begin{tikzpicture}[
my legend style/.style={
	rectangle,
	draw,
	very thin,
	minimum width=15pt,
},
laboratory style/.style={
	fill=plots-laboratory,
	mark=mysquare,
	every mark/.append style={fill=plots-laboratory},
	every path/.style={postaction={
			nomorepostaction,
			pattern=north east lines,
			pattern color=white,
	}}
},
laboratory legend style/.style={
	my legend style,
	preaction={fill, plots-laboratory},
	pattern=north east lines,
	pattern color=white
},
online style/.style={
	fill=plots-online,
	mark=mysquare,
	every mark/.append style={fill=plots-online},
	every path/.style={postaction={
			nomorepostaction,
			pattern=north west lines,
			pattern color=white,
	}},
},
online legend style/.style={
	my legend style,
	preaction={fill, plots-online},
	pattern=north west lines,
	pattern color=white
},
mturk style/.style={
	fill=plots-mturk,
	mark=mysquare,
	every mark/.append style={fill=plots-mturk},
	every path/.style={postaction={
			nomorepostaction,
			pattern=crosshatch dots,
			pattern color=white,
	}},
},
mturk legend style/.style={
	my legend style,
	preaction={fill, plots-mturk},
	pattern=crosshatch dots,
	pattern color=white
}
]
\begin{axis}[name=plot-main,
	scale only axis,
	hide y axis,
	axis x line*=bottom,
	xmin=0,
	xmax=105,
	ymin=-1.5,
	ymax=14.5,
	y=7pt,
	mark size=4pt,
	width=0.7\linewidth-1\axislabellength,
	my boxplot/.style={
		boxplot={
			box extend=0.8,
			whisker extend=0.6
		},
		fill,
		draw=black,
		very thin,solid,
	},
	]
	\participantsplots{}
\end{axis}

\begin{axis}[name=plot-outlier,
	at={($(plot-main.east)+(-0.0cm,0cm)$)},anchor=west,
	scale only axis,
	hide y axis,
	axis x line*=bottom,
	axis x discontinuity=parallel,
	xmin=110,
	xmax=430,
	ymin=-1.5,
	ymax=14.5,
	y=7pt,
	mark size=4pt,
	width=0.3\linewidth-0*\axislabellength,
	my boxplot/.style={
		boxplot={
			box extend=0.7,
			whisker extend=0.6
		},
		fill,
		draw=black,
		very thin,solid,
	},
	nodes near coords={\pgfmathprintnumber{\pgfkeysvalueof{/data point/x}}},
	every node near coord/.style={
		color=black,
		font=\footnotesize,
	},
	]
	\participantsplots{}
	
	\node[anchor=east] at (axis cs:430,13) { VTC };
	\node[anchor=east] at (axis cs:430,9) { Configuration };
	\node[anchor=east] at (axis cs:430,5) { Usability };
	\node[anchor=east] at (axis cs:430,1) { Meta };
\end{axis}

\path[] (plot-main.south) |- (0.5\linewidth, 0) node[pos=1.0] (legend-anchor) {};
\matrix[below=3pt of legend-anchor,column sep=0.5cm,row sep=0.5cm,anchor=north]{
	\node[laboratory legend style,label={[text depth=0.0pt]0:Laboratory}] {}; &
	\node[online legend style,label={[text depth=0.0pt]0:Online}] {}; &
	\node[mturk legend style,label={[text depth=0.0pt]0:Mechanical Turk}] {};\\
};
\end{tikzpicture}%%
	\vspace{-1.2em}
	\caption{Number of participants categorized by evaluation practice and study implementation.}%
	\label{fig:participants}%
\end{figure*}%

\subsection{Summary of Findings}
\label{sec:summary:findings}
The results of the studies are thematically widely scattered and range from VTC to eye-movement exploration patterns.
To extract and summarize these results, statistically significant findings of the studies were reviewed and verified instead of the respective claims.
These summarized findings are not intended for citation as they are broadly reduced statements with a too general meaning---without further revalidation and replication, we insinuate that the measured effects are not generalizable by the studies' designs; refer to the original publications instead.
The dependent variables measures are stated as follows: effectiveness (E), efficiency (I), and user satisfaction (S)~\cite{leroy2011-designing}.



\begin{table}[t]
  \centering
  \scriptsize%
  \noindent%
  \renewcommand{\arraystretch}{1.01}
  \rowcolors{3}{white}{gray!20}%
  \setlength{\tabcolsep}{1.7pt}%
  \setlength{\aboverulesep}{0pt}
  \setlength{\belowrulesep}{0pt}
  \setlength{\arrayrulewidth}{0.4pt}
  \caption{Visualization techniques that were compared to treemaps in VTC studies.
  Notation: evaluation with (\significOK) and without (\notsignificOK) significant results.\\ Abbreviation: State of the Practice \emph{(StotP)}.\vspace{-0.5em}}%
    \begin{tabular}{                                                    r|c!{\tls}     c!{\tls}      c!{\tls}          c                |c!{\tls}         c!{\tls}         c!{\tls}            c!{\tls}         c                     |c!{\tls}              c!{\tls}     c!{\tls}                 c               |c            |}
                                                                         & \mc{4}{traditional}                                          & \mc{5}{recent}                                                                              & \mc{4}{StotP}                                                              &             \\
                                                                         & \rot{\strut List} & \rot{\strut Table} & \rot{\strut Bar Chart} & \rot{\strut Pie Chart} & \rot{\strut Nodelink} & \rot{\strut Sunburst} & \rot{\strut Icicle Plot} & \rot{\strut Beamtree} & \rot{\strut Botanical Tree} & \rot{\strut Windows Expl.} & \rot{\strut UNIX} & \rot{\strut Eclipse \& Excel} & \rot{\strut Printout} & \rot{\strut Other} \\\hline
      Directory Browsing~\cite{johnson1993-categorical, turo1992-improve} &            &             &                 &                 &                &                &                   &                &                      &                     & \significOK        &                        &                &             \\
                                 GEnie Experiment~\cite{turo1992-improve} &            &             &                 &                 &                &                &                   &                &                      &                     &            &                        & \significOK            &             \\
                      US Budget Experiment~\cite{johnson1993-categorical} &            & \significOK         &                 &                 &                &                &                   &                &                      &                     &            &                        &                &             \\
                3D Information Visualizations~\cite{wiss1999-tasksupport} &            &             &                 &                 & \significOK            &                &                   &                &                      &                     &            &                        &                &             \\
                Sunburst vs. Treemap (small)~\cite{stasko2000-structures} &            &             &                 &                 &                & \significOK            &                   &                &                      &                     &            &                        &                &             \\
                Sunburst vs. Treemap (large)~\cite{stasko2000-structures} &            &             &                 &                 &                & \significOK            &                   &                &                      &                     &            &                        &                &             \\
                       Topology \& Node Size~\cite{barlow2001-comparison} &            &             &                 &                 & \significOK            & \significOK            & \significOK               &                &                      &                     &            &                        &                &             \\
                         Compare to Beamtrees~\cite{vanHam2003-beamtrees} &            &             &                 &                 &                &                &                   & \significOK            &                      &                     &            &                        &                &             \\
                   Compare 5 Tree Visualizations~\cite{kobsa2004-systems} &            &             &                 &                 & \significOK            &                &                   & \significOK            & \significOK                  & \significOK                 &            &                        &                &             \\
                        Treemaps \& Tables~\cite{goldberg2005-monitoring} &            & \significOK         &                 &                 &                &                &                   &                &                      &                     &            &                        &                &             \\
                                 Crop Circles~\cite{wang2006-cropcircles} &            &             &                 &                 & \significOK            &                &                   &                &                      &                     &            &                        &                &             \\
                   Knowledge Discovery~\cite{wang2006-knowledgediscovery} &            &             &                 &                 & \notsignificOK       & \significOK            &                   &                &                      &                     &            &                        &                &             \\
                                Hierarchy Browser~\cite{andrews2007-hvte} &            &             &                 &                 & \significOK            & \significOK            &                   &                &                      &                     &            &                        &                &             \\
                        Search Engine Results~\cite{chu2007-searchengine} & \significOK        &             &                 &                 &                &                &                   &                &                      &                     &            &                        &                &             \\
                             Multiple Views~\cite{teoh2007-multipleviews} & \significOK        &             &                 &                 &                &                &                   &                &                      &                     &            &                        &                &             \\
                 Search Result Interface 1~\cite{clarkson2009-resultmaps} & \significOK        &             &                 &                 &                &                &                   &                &                      &                     &            &                        &                &             \\
                 Search Result Interface 2~\cite{clarkson2009-resultmaps} & \significOK        &             &                 &                 &                &                &                   &                &                      &                     &            &                        &                &             \\
                                  Data Density~\cite{kong2010-guidelines} &            &             & \significOK             &                 &                &                &                   &                &                      &                     &            &                        &                &             \\
                            Node Link Treemaps~\cite{linsen2011-nodelink} &            &             &                 &                 & \notsignificOK &                &                   &                & \notsignificOK & \notsignificOK            &            &                        &                &             \\
              Software Systems as Cities~\cite{wettel2011-softwarecities} &            &             &                 &                 &                &                &                   &                &                      &                     &            & \significOK                    &                &             \\
                Eye-tracking for EER-Tutor~\cite{mathews2012-eyetracking} &            &             &                 &                 & \notsignificOK       &                &                   &                &                      &                     &            &                        &                & \notsignificOK    \\
                            Rapid Overview~\cite{mueller2013-eyetracking} &            &             &                 &                 & \significOK            &                & \significOK               &                &                      &                     &            &                        &                &             \\
                 Software Landscape~\cite{fittkau2015-softwarelandscapes} &            &             &                 &                 & \significOK            &                &                   &                &                      &                     &            &                        &                &             \\
                      Treemap vs. Euler~\cite{sathiyanarayanan2015-euler} &            &             &                 &                 &                &                &                   &                &                      &                     &            &                        &                & \notsignificOK    \\
                    Treemap vs. Euler 2~\cite{sathiyanarayanan2015-graph} &            &             &                 &                 &                &                &                   &                &                      &                     &            &                        &                & \notsignificOK \\
                              Photo: Layout~\cite{guerra-gomez2016-photo} & \significOK        &             &                 &                 &                &                &                   &                &                      &                     &            &                        &                &             \\
              Compare 3 Tree Visualizations~\cite{mueller2017-perception} &            &             &                 &                 & \significOK            & \significOK            &                   &                &                      &                     &            &                        &                &             \\
                   Wrapped Bars: Comparison~\cite{yalcin2017-wrappedbars} &            &             & \significOK             &                 &                &                &                   &                &                      &                     &            &                        &                &             \\
                      Wrapped Bars: Ranking~\cite{yalcin2017-wrappedbars} &            &             & \significOK             &                 &                &                &                   &                &                      &                     &            &                        &                &             \\
                 Wrapped Bars: Distribution~\cite{yalcin2017-wrappedbars} &            &             & \significOK             &                 &                &                &                   &                &                      &                     &            &                        &                &             \\
                    Pie Charts: Proportion~\cite{huestegge2018-processes} &            &             &                 & \significOK             &                &                &                   &                &                      &                     &            &                        &                &             \\
                    Pie Charts: Comparison~\cite{huestegge2018-processes} &            &             &                 & \significOK             &                &                &                   &                &                      &                     &            &                        &                &             \\
                      Pie Charts: Labeling~\cite{huestegge2018-processes} &            &             &                 & \significOK             &                &                &                   &                &                      &                     &            &                        &                &             \\
                Pie Charts \& Stacked Bars~\cite{huestegge2018-processes} &            &             &                 & \significOK             &                &                &                   &                &                      &                     &            &                        &                &             \\\hline
    \end{tabular}%
    \vspace{-1em}
    \label{tab:vtc}%
\end{table}%


\paragraph{\textbf{Visualization Technique Comparison Studies}}
In comparison to recent visualization techniques, treemap performance results are heterogeneous.
\emph{Sunburst views}~\cite{stasko2000-structures,barlow2001-comparison} and \emph{icicle plots}~\cite{barlow2001-comparison} perform better\textsuperscript{EIS} than treemaps for tasks that refer to tree attributes.
Likewise, \emph{icicle plots} perform better\textsuperscript{E} than treemaps for tasks related to counting and size~\cite{mueller2013-eyetracking,mueller2017-perception}.
Titlebar treemaps and cushion treemaps perform better\textsuperscript{EIS} than \emph{3D beamtrees} and \emph{botanical trees}~\cite{kobsa2004-systems}---except for tasks related to level counting where 3D beamtrees perform better\textsuperscript{EI}~\cite{vanHam2003-beamtrees,kobsa2004-systems}.
Findings targeting treemaps and \emph{node-link diagrams} are diversified and partly indecisive.
Traditional data visualization techniques were in most cases worse than treemaps when used on non-tree-structured data.
Treemaps are more efficient than \emph{tables}~\cite{johnson1993-categorical,goldberg2005-monitoring} and have a higher accuracy than \emph{pie charts}~\cite{huestegge2018-processes} for comparison tasks.
For dense non-hierarchical sorted numerical data~\cite{yalcin2017-wrappedbars}, treemaps perform worse\textsuperscript{E}.
On the other hand, treemaps perform better\textsuperscript{E} than \emph{bar charts}~\cite{kong2010-guidelines} for both high-density two-level hierarchies and comparison of non-leaf nodes.
For search results displayed in a \emph{list}, an additional treemap visualization improves the user's satisfaction and understanding of the results~\cite{chu2007-searchengine,clarkson2009-resultmaps}.

Treemaps outperform other state-of-the-art exploration tools in terms of efficiency.
Cushion treemaps, for example, are rated lower than \emph{Windows Explorer}~\cite{kobsa2004-systems,andrews2007-hvte} and perform worse\textsuperscript{EI} for tasks that refer to the complete (sub-)tree~\cite{kobsa2004-systems}.
Then again, treemaps are more efficient and are rated better for these tasks than the \emph{UNIX shell}~\cite{johnson1993-categorical,turo1992-improve}.
Nested treemaps are more efficient for structure-similarity tasks than Windows Explorer~\cite{wang2006-knowledgediscovery}.
For program comprehension tasks, treemaps perform better\textsuperscript{EI} than \emph{Eclipse \& Excel}, while being especially more efficient for leaf-node extrema tasks~\cite{wettel2011-softwarecities}.
Likewise, treemaps are more efficient for global-scope tasks than \emph{printout} financial reports~\cite{turo1992-improve}.

\paragraph{\textbf{Design Element Studies}}
Multiple user studies examined new and existing \emph{visual variables} in the context of treemaps.
\emph{Luminance} was found to be suitable to encode secondary attributes when size encodes primary values~\cite{kong2010-guidelines} and for \emph{color} the C-VPA coloring method showed improved size discrimination and higher aesthetic rankings than random coloring method~\cite{xie2017-coloraesthetics}.
Newly proposed visual variables are \emph{sketchiness} in 2.5D~\cite{limberger2016-sketchy} and \emph{rectangular orientation} in 2D~\cite{liang2012-orientation}, both can be used as independent visual variables without noticeable impact on other visual variables.
It should be noted that the accuracy of 2D \emph{size} judgement tasks suffers when the compared rectangles have extreme \emph{aspect ratio}s or when the rectangles are squares~\cite{kong2010-guidelines}.
The findings of studies that compare \emph{layout algorithms} are so extensive that it is difficult to summarize them in this report and are therefore future work~\cite{bederson2002-quantum,hahn2017-directionchange,zhou2017-cartography}.
\par
The \emph{depiction of topology} influences the perception of node size and structure~\cite{johnson1993-categorical}.
For example, titlebar treemaps perform better\textsuperscript{EI} than cushion treemaps which perform better than non-nested treemaps for tasks that refer to structure~\cite{kobsa2004-systems,irani2006-perception,biuk-aghai2017-maplike}.
Likewise, titlebar treemaps perform better\textsuperscript{IS} than simple nested treemaps for identification tasks and comparison tasks~\cite{goldberg2005-monitoring}.
In addition to the kind of topology depiction, the \emph{shape} changes the perception of structure and interpretation of relations.
For relationship identification tasks, rectangular treemaps are less effective but are more efficient than map-like treemaps~\cite{biuk-aghai2017-maplike}.
However, circular treemaps might have a better performance for structure extrema tasks than rectangular treemaps~\cite{wang2006-cropcircles}.
For different \emph{dimensionalities}---especially 2D and 2.5D---differences could not be confirmed~\cite{bladh2004-threedimensions}.
\par
Suitable \emph{focus and context} techniques for treemaps are Balloon Focus~\cite{tu2008-multifocus} and Drill-down~\cite{shi2005-contentbrowsing}.
Users of animated treemaps will take more likely shortcuts and rate them higher than users of non-animated treemaps~\cite{bladh2005-transitions}.
As a supporting technique, node \emph{aggregation} is recognized by users and tasks that refer to determining nodes of interest are rated easier than without aggregation~\cite{limberger2017-aggregation}.
Similarly, in 2.5D treemaps, interactive \emph{height-based filtering} might increase the accuracy of basic height-related tasks with only a marginal interaction overhead~\cite{limberger2018-filtering}.
In \emph{image treemaps}, users prefer fewer images and the presence of labels~\cite{guerra-gomez2016-photo}, while two-level-clustering reduces image finding time~\cite{gomi2008-imagebrowser}.

\paragraph{\textbf{Usability Studies}}
During the exploration of a data set in a treemap visualization, certain \emph{visual exploration strategies} apply.
Treemaps are effective in promoting task-relevant fixations~\cite{mueller2017-perception} and improve the user's perception of some patterns~\cite{wang2006-knowledgediscovery}.
To explore these patterns in the underlying data treemap users used node, node count, node type, node size, structure, and content information~\cite{wang2006-knowledgediscovery}.
However, treemap users constantly have to reorient and remember which elements were already processed~\cite{mueller2017-perception}.
The simultaneous usage of complementary visualization techniques may increase the user's understanding of the visualized data~\cite{teoh2007-multipleviews}.
By their way of visualizing structure, treemaps are more or less suited for certain \emph{task types}.
In contrast to size-related tasks where treemaps offer a good sense of this attribute even between nodes of different levels~\cite{teoh2007-multipleviews}, tasks related to global structure and some detailed structure-related tasks---level and sibling tasks---seem to be difficult to solve with treemaps~\cite{kobsa2004-systems}.
\par
Remaining findings include \emph{shape}, \emph{dimension}, and \emph{interaction}.
In map-like treemaps, users have a strong location recall with a clear understanding of islands, whereas landmarks might hinder interpretation of relationship discontiguity~\cite{hawes2015-codesurveyor}.
While in \emph{virtual reality} the interaction with 2.5D treemaps can be implemented by rotation, translation, and selection gestures~\cite{fittkau2015-virtualreality}; 3D treemaps are %
suitable for exploration and long-term learning tasks~\cite{wiss1999-tasksupport}.

\paragraph{\textbf{Meta Studies}}
The task solving performance of visualization users is influenced by the compatibility of visual and \emph{verbal metaphor}. 
This complex effect is varied by the user's preference of metaphor, spatial ability, and the locus of control~\cite{ziemkiewicz2008-metaphors,ziemkiewicz2011-locus}.
Beyond that, the visualization technique with the highest \emph{visual quality} seems to perform relatively high in effectiveness-related measures~\cite{cawthon2007-aesthetic}.
\paragraph{\textbf{Unconsidered Studies}}
Some user studies do not meet the requirements to be considered for further analysis.
Due to a lack of a clear description of the results, the user studies on Euler diagrams were not considered~\cite{sathiyanarayanan2015-euler, sathiyanarayanan2015-graph}.
Second, we did not consider a study targeting the combination of treemaps and node-link diagrams since the methodical approach is not appropriate~\cite{linsen2011-nodelink}.
The authors tried to extend a previous user study~\cite{kobsa2004-systems} to compare their results with the other evaluation results.
Due to a lack of comparability to the former user study---the used data set was different and the statistical data was not available---and no other comparable visualization that were evaluated within the same user study, the results were not considered.
Similarly, we had to discard a user study using eye tracking as the visualized amount of data differed between the compared visualizations~\cite{mathews2012-eyetracking}.
Last, we do not consider a user study targeting photo arrangements using a treemap approach~\cite{guerra-gomez2016-photo}, as the complex interactions and learning effects were not discussed.

\section{Critical Review} %
Following, the treemap user studies are reviewed according to their quality, meaningfulness, and the resulting usefulness for the research community.
The objective is to determine whether the findings can be accepted as valid and, thus, considered as generalized statements.
If not, the main problems and shortcomings encountered in the studies are briefly analyzed.

\subsection{Reviewing the User Studies}
In order to assess the quality of these user studies individually as well as collectively, we reviewed the comprehensibility, fairness, and completeness of particular studies and the topic coverage across all studies.
On the basis of the research goal and the hypotheses derived, we have to conclude that many studies lack this information.
Half of the user studies do not state their hypotheses based on a precise research goal, although they are the basis of every study. 
Several texts fail to comprehensively describe the study design and present their variable design.
Almost no study discusses the interfering factors and the concluding validity issues. 
In addition to these design issues, the variable conditions in controlled experiments are often not randomized.
Other frequently missing information includes detailed descriptions of data sets, participants, and procedure as well as values regarding statistical significance.
About a quarter of the studies do not report on their participants' background.
Typically, participants of treemap studies have none of domain knowledge, previous knowledge of treemaps, or experience therewith.
Only a small amount of studies invited domain experts.
The most common participant type are students of computer science or human-computer interaction.
The number of participants is often too small to provide the statistical power or to serve the number of independent variables, as reported in previous surveys.
Likewise, claims---especially stated in abstracts and conclusions---go far beyond the results.
Even for studies that are not controlled experiments, results are sometimes verbally stated as valid generalized statements based on a causation. 
\par
Although comprehensive evaluation intends multiple points in time for assessment, the most common point is at or near the end of the development process.
At this point, many researchers compare their approach with previous ones and try to prove that their approach is equally good at least, which requires a controlled experiment.
This reduces the versatility of the study types used and thus the scope of the evaluation findings.
Typically, such studies also test features of the user interface that are often not evaluated separately from the underlying visualization technique.
Often, user studies derive their usage scenarios and tasks from previous publications, although no theoretical background is available.
Since there is no standardized data set, each study uses own data set.
Despite three consecutive studies by Ziemkiewicz~et~al.~\cite{ziemkiewicz2008-metaphors,ziemkiewicz2009-preconceptions,ziemkiewicz2011-locus}, there has not been one dedicated replication study on treemaps so far.
\paragraph{\textbf{Task Compatibility}} %
With the demand for fair control conditions, the configuration of a treemap depends on the task~\cite{johnson1993-categorical}.
With the exception of Johnson's doctoral thesis~\cite{johnson1993-categorical} and configuration studies, most publications do not discuss whether the chosen configuration is appropriate.
For tasks that refer to structure, treemaps with a larger offset or titlebar treemaps are appropriate, whereas size-related tasks require no offset~\cite{johnson1993-categorical}.
Treemaps with cushion topology depiction~\cite{vanwijk1999-cushion} offer a trade-off between both needs.
About 40\,\% of reviewed user studies contain tasks that refer to structure or tasks that refer to tree attributes, which are based on structure.
Most of these studies configured their treemaps accordingly with offset or cushion.
Whereas for the most common task attribute \emph{size}, only one third of the studies are properly configured. %
The frequently used counting tasks---which appears in almost one fifth of the studies---in general do not fit visualization techniques that do have no support for a user interface.
According to Ziemkiewicz, the matching of visual and verbal metaphor is important~\cite{ziemkiewicz2008-metaphors}.
Although treemaps visually rely on the containment metaphor, about half of the studies used the level metaphor in their task descriptions.
Treemaps are able to simultaneously show multiple visual variables and the structure of a data set.
Nevertheless, only one third of the studies included tasks that refer to two visual variables.
Only four studies contain tasks targeting three dimensions, and correlation studies are the least common task type.

\paragraph{\textbf{Coverage of Evaluation Topics}}
In order to assess the user studies as an entirety, we use the proposed evaluation topics to rate coverage and identify underrepresented areas.
With respect to \emph{Use Cases}, mainly VTC studies are relevant.
Eleven techniques were compared and the results refer to three tasks types suitable for treemaps.
An additional topic was the simultaneous visualization of treemaps and other techniques.
While studies related to use cases include multiple studies, there is only one topic covered for \emph{Communication}. 
With regard to \emph{Interaction} aspects, the reviewed user studies cover pattern matching and multiple focus-and-context techniques including filtering in 2.5D, node aggregation, and animation.
The connection of visual quality and the encouragement of users was also evaluated.
The entirety of user studies considers several \emph{Configurations}.
From our assessment, questions regarding dimensionality are still open.
The studies offer recommendations for shape, suitable to two tasks types extended by few qualitative insights.
For the selection of an appropriate topology depiction, there are three recommendations related to different task types.
In total, six visual variables were examined, whereby two were new in the context of treemaps.
Treemap layout studies require a systematic review including non-participant measures~\cite{vernier2020-comparison}.
Additionally, \emph{Perceptional} aspects covered are task relevant fixations, area perception, and problems in orientation.
Moreover, treemaps seem to lack a natural understanding of their basic containment metaphor without adequate training, such as the size of inner nodes~\cite{mueller2017-perception}.
In summary, treemap user studies suggested different appropriateness of treemaps in comparison to other visualization techniques related to six task types. 
While visual variables are well covered, there is hardly any validation of data set topology.
The topics \emph{Perception} and \emph{Communication} are underrepresented.
For treemap configuration, studies offer some suggestions regarding topology depiction, while dimensionality remains an open question.
Further, treemap user studies seem to cover multiple aspects of treemaps but are fairly unstructured.
Multiple topics are underrepresented, while others convey little structure.
\subsection{Deriving Generalized Statements}%
\label{sec:discussion}
Based on our review, we conclude that there is little that can be clearly derived.
Various shortcomings and quality issues in study design, preparation, and execution were found.
Other points to be added are the conclusiveness of research objective and study design as well as the analysis of nuisance or confounded variables, multiple participant related issues referring to count and representativeness of sample, and the lack of formal documentation of setup and conduct.
In addition, almost no study addresses these problems in a validity discussion.
There are also no replication studies.
Concluding, this hinders general statements about treemaps, since most of the studies are affected.
The presented study results should be used as hints and guides for future research and motivate replication studies of these.
Researchers should thoroughly review previous studies before referring to their results or claims.

There is still little knowledge on which task types---and thus use cases or usage scenarios---are truly suitable to treemaps.
Despite most-basic, theoretical background, there are no recommendations or guidelines on configuration that visualization designer could access.
Likewise, many perceptional aspects of treemaps are unclear.
For example, there is no resilient data on the value of simultaneous usage of more than one visual variable---although it is frequently referred to.
Advantages or disadvantages of 2.5D treemaps or 3D treemaps are still unclear.
Further, open questions are suitable ranges of data set sizes, learnability of the containment metaphor, and the mental-map characteristic of treemaps referring to long-term understanding of data.
With regards to a holistic view on treemaps, the previous studies are historically grown and lack an inherent structure.
Up to date we know of no research agenda to structurally prove previous statements about treemaps and evaluate the remaining aspects.
This might be caused by missing references for a visualization evaluation standard, although there are initial thoughts on that topic~\cite{zuk2006-heuristics}.

\section{Conclusion}
\label{sec:conclusions}
Although a large number of user studies on treemap visualization techniques were designed, conducted, and published, there is still little that can be reliably derived or generalized.
We addressed this gap by means of an analysis and review of the current state of research as well as a prospect for future researchers.
We provided a summary, excerpt and review of 69 user studies that are additionally explorable on a community website.
Although this list of studies is far from being complete, these resources can be used as a starting point for researchers to ease design, execution, and evaluation of user studies that target treemaps.
We believe this evaluation, the review of previously published user studies, and the summary of guidelines serve as a starting point for future user studies and structure the evaluation process, leading to further insights into treemap visualization techniques.
\par
As part of future work, we plan to prepare an overview of guidelines and conceptual models on user studies in information visualization to provide a systematic approach for researchers.
A next step would be a research agenda with separate directions and projected effects.
The provided data set is intended as a starting point and is target to get extended by the authors themselves to include evaluations of recently and future conducted user studies.
Likewise, the insights from user studies targeting other tree and hierarchy visualization techniques are another area of interest for integration into the data set.
But most importantly, we invite the information visualization community to collaborate on the user study data set.

\begin{acks}
We want to thank the anonymous reviewers for their valuable comments and suggestions to improve this article.
This work has been supported by the German Federal Ministry of Education and Research (BMBF) through grant 01IS19006  (\enquote{KI-Labor ITSE}).
Further, this work is part of the \enquote{Software-DNA} project, which is funded by the
European Regional Development Fund (ERDF -- or EFRE in German) and the
State of Brandenburg (ILB) as well as the \enquote{TASAM} project, which is funded by the German Federal Ministry for Economic Affairs and Energy (BMWi, ZIM).
\end{acks}

\bibliographystyle{ACM-Reference-Format}
\bibliography{references,referencesStudies}

\end{document}
\endinput

